\chapter{Implementare}
\label{ch:implementare}

\section{Motorul \texttt{memtable}}

Operația de citire dintr-un shard ilustrează politicile de evacuare: expirarea TTL este verificată leneș
la acces, iar efectul unui hit depinde de politică.

\begin{code}{Citirea dintr-un shard (\texttt{crates/memtable/src/shard.rs}).}
pub(crate) fn get(&mut self, key: &str, now: Instant) -> Lookup {
    let Some(&i) = self.map.get(key) else { return Lookup::Miss };
    if self.entry(i).expires_at.is_some_and(|e| e <= now) {
        self.remove(i);              // expirare leneșă
        return Lookup::Expired;
    }
    self.touch(i);                   // LRU: mută în cap; CLOCK: bit = 1; FIFO: nimic
    Lookup::Hit(self.entry(i).value.clone())   // Bytes: clonare O(1), fără copiere
}
\end{code}

Evacuarea CLOCK parcurge coada: intrările cu bitul de referință setat primesc a doua șansă. Bucla se
termină după cel mult o trecere completă, deoarece fiecare intrare vizitată își pierde bitul.

\begin{code}{Evacuarea unei intrări.}
fn evict_one(&mut self) -> bool {
    loop {
        let t = self.tail;
        if t == NIL { return false; }
        if self.policy == Policy::Clock && self.entry(t).referenced {
            self.entry_mut(t).referenced = false;   // a doua șansă
            self.unlink(t);
            self.push_front(t);
            continue;
        }
        self.remove(t);
        return true;
    }
}
\end{code}

Statisticile sunt contoare atomice (\texttt{AtomicU64} cu ordonare \texttt{Relaxed}), actualizate în
afara secțiunii critice.

\section{Nodul de cache și clientul}

Serverul folosește un task tokio per conexiune. Bucla de tratare citește o linie de antet, apoi exact
\texttt{len} octeți de payload, și scrie răspunsul într-un \texttt{BufWriter}; flush-ul se face doar
când nu mai există comenzi în buffer. Erorile de protocol sunt raportate clientului fără a închide
conexiunea; erorile de I/O închid conexiunea. Un janitor elimină periodic intrările expirate.

Clientul păstrează un pool de conexiuni (\texttt{Mutex<Vec<BufStream<TcpStream>>>}); mutex-ul este
ținut doar pentru \texttt{pop}/\texttt{push}, niciodată în timpul unui \texttt{await}. Cheile și tag-urile
sunt validate la granița de încredere (fără spații sau virgule), iar valorile sunt limitate la 16 MiB pe
server.

\section{Aplicația web}

\subsection{Backend-uri interschimbabile}

Backend-urile sunt modelate ca \texttt{enum}, cu \emph{dispatch} static:

\begin{code}{Backend-urile de cache (\texttt{crates/webapp/src/cache.rs}).}
pub enum Backend {
    None,
    Local(Arc<Table>),                          // L1 în proces
    Cluster(ClusterClient),                     // noduri cache-node
    Redis(RedisConn),
    PgMem(PgPool),                              // tabelă UNLOGGED pe tmpfs
    Tiered { l1: Arc<Table>, l2: ClusterClient }, // near-cache
}
\end{code}

Configurația activă (\texttt{Runtime}: backend + parametri) este ținută într-un
\texttt{RwLock<Arc<Runtime>>}. Fiecare cerere clonează doar \texttt{Arc}-ul (o incrementare atomică), iar
schimbarea modului înlocuiește atomic pointerul — cererile în curs își termină execuția pe vechea
configurație.

\subsection{Single-flight}

\begin{code}{Single-flight cu \texttt{OnceCell} (\texttt{state.rs}, simplificat).}
let (cell, leader) = {
    let mut map = self.inflight.lock();
    match map.get(key) {
        Some(c) => (Arc::clone(c), false),
        None => { let c = Arc::default(); map.insert(key.into(), Arc::clone(&c)); (c, true) }
    }
};
let res = cell.get_or_init(move || async move { fetch.await }).await.clone();
// Oricine termină curăță intrarea (leader-ul poate fi anulat de deconectarea clientului).
let mut map = self.inflight.lock();
if map.get(key).is_some_and(|c| Arc::ptr_eq(c, &cell)) { map.remove(key); }
\end{code}

Un detaliu important: dacă cererea „lider" este anulată (clientul închide conexiunea), \texttt{OnceCell}
permite altei cereri să preia inițializarea, iar curățarea intrării este făcută de oricine termină —
altfel cheia ar rămâne blocată pe un rezultat vechi.

\subsection{API, documentație și interfață}

Rutele sunt definite cu axum 0.8; documentația OpenAPI 3 este generată din cod cu \texttt{utoipa}
(adnotări \texttt{\#[utoipa::path]} pe handlere și \texttt{ToSchema} pe tipuri) și servită prin Swagger UI
la \texttt{/swagger-ui}. Răspunsurile includ antetul \texttt{X-Cache}. Interfața web (\texttt{/}) este o
pagină statică inclusă în binar care permite schimbarea modului pe toate instanțele, afișează live rata de
hit și operațiile pe secundă (Chart.js), conține un generator de trafic Zipf rulat în browser și un
\emph{playground} pentru API.

\subsection{Observabilitate}

Un middleware măsoară latența fiecărei cereri într-o histogramă Prometheus etichetată cu ruta
(\texttt{MatchedPath}), metoda și codul de stare. Alte metrici: lookup-uri hit/miss, interogări DB pe tip,
erori ale backend-ului de cache și adâncimea cozii write-behind. Grafana este provizionată automat cu
sursa de date și un dashboard cu șase panouri.

\section{Optimizări}
\label{sec:optimizari}

Optimizările au fost ghidate de măsurători: o rulare scurtă (\texttt{scripts/profile.sh}) cu snapshot
de consum CPU per container, iterată după fiecare modificare.
\ifdefined\OptTable\OptTable\fi

\paragraph{1. Profilare: baza de date era gâtul de sticlă.} Înainte de optimizare, PostgreSQL consuma
$\approx$410\% din cele 500\% CPU disponibile, în timp ce fiecare instanță web consuma sub 20\%. Planul de
execuție al clasamentului unei categorii arăta un \emph{Parallel Seq Scan} pe 100.000 de produse și 500 de
accesări de heap pe \texttt{reviews} (7{,}6~ms).

\paragraph{2. Tuning PostgreSQL.} (a) \texttt{max\_parallel\_workers\_per\_gather = 0}: pentru un
volum OLTP mare, pornirea de workeri paraleli pentru fiecare interogare costă mai mult decât câștigă;
(b) indecși acoperitori \texttt{reviews(product\_id) INCLUDE (rating)} și
\texttt{products(category\_id) INCLUDE (id, name, price\_cents)} plus \texttt{VACUUM ANALYZE}, care
permit \emph{index-only scans}. Interogarea a scăzut de la 7{,}6~ms la 1{,}5~ms.

\paragraph{3. Calea fierbinte din aplicație.} Eliminarea \texttt{TraceLayer} (un span de tracing per
cerere), contoare Prometheus pre-rezolvate (fără căutarea etichetelor la fiecare cerere), evitarea
alocărilor în middleware, allocatorul \texttt{mimalloc} și reducerea pool-ului DB de la 32 la 16
conexiuni per instanță (mai puțină contenție în PostgreSQL).

\paragraph{4. \texttt{NOTIFY} condiționat.} Eșantionarea \texttt{pg\_stat\_activity} sub încărcare a arătat
majoritatea sesiunilor de scriere blocate în \texttt{Lock:object}: o tranzacție care execută
\texttt{NOTIFY} ține un lock global pe coada de notificări până după flush-ul WAL la commit, ceea ce
serializează toate scrierile și anulează \emph{group commit}. Trigger-ul notifică acum doar dacă sesiunea
nu a setat \texttt{app.notify = off}; aplicația setează valoarea (local tranzacției, prin
\texttt{set\_config(\ldots, true)} în aceeași instrucțiune \texttt{UPDATE}) când backend-ul nu are L1.
Scrierile externe continuă să fie notificate.

\paragraph{5. Agregate materializate.} Chiar și după pașii anteriori PostgreSQL rămânea la
$\approx$260\% CPU. Cauza: fiecare modificare de preț invalidează clasamentul categoriei; cu
$\approx$2.200 de scrieri/s distribuite pe 200 de categorii, aproape orice cerere de clasament era un miss
și recalcula un \texttt{JOIN} + \texttt{GROUP BY} de 1{,}5~ms. Media notelor și numărul de recenzii au
devenit coloane în \texttt{products} (\texttt{rating\_avg}, \texttt{review\_count}), întreținute incremental
de un trigger pe \texttt{reviews}, iar clasamentul folosește indexul
\texttt{(category\_id, rating\_avg DESC, review\_count DESC, id)}: interogarea durează acum
$\approx$0{,}09~ms (de 85 de ori mai puțin decât forma inițială). Citirea unui produs devine un lookup pe
cheia primară, iar scrierea întoarce rândul actualizat prin \texttt{RETURNING} (o interogare în loc de două).
Prețul nu este indexat și tabela are \texttt{fillfactor = 90}, deci actualizările de preț rămân
\emph{HOT} (nu modifică indecșii).

\paragraph{Rezultat.} Pe aceeași mașină, debitul cu cache local a crescut de la
\R{opt/baseline/rps} la \R{opt/mat-local/rps} cereri/s, iar p99 a scăzut de la \R{opt/baseline/p99}~ms la
\R{opt/mat-local/p99}~ms; varianta fără cache a crescut de la \R{opt/baseline-none/rps} la
\R{opt/mat-none/rps} cereri/s. Cu cache, consumul PostgreSQL a scăzut de la $\approx$410\% la
$\approx$60–100\% CPU; cei mai mari consumatori rămân nginx și generatorul de încărcare. Toate experimentele
din \cref{ch:evaluare} folosesc versiunea finală.

\section{Calitatea codului}

\begin{itemize}
  \item \textbf{Lint}: \texttt{clippy} cu grupurile \texttt{all} și \texttt{pedantic} la nivel de workspace,
        rulat cu \texttt{-D warnings}; \texttt{unsafe\_code = "forbid"}.
  \item \textbf{Formatare}: \texttt{rustfmt} (configurat în \texttt{rustfmt.toml}), \texttt{.editorconfig}.
  \item \textbf{Dependențe}: \texttt{cargo-deny} (licențe, vulnerabilități, surse) prin \texttt{deny.toml}.
  \item \textbf{Teste}: 19 teste unitare și de integrare (politici de evacuare, TTL, tag-uri,
        concurență, echilibrul inelului, roundtrip client–server, cluster, simulator) și un test end-to-end
        de coerență (\texttt{scripts/consistency.sh}) pentru toate cele 18 combinații backend × strategie.
  \item \textbf{CI}: GitHub Actions rulează formatarea, clippy, testele, \texttt{cargo-deny} și
        construirea imaginii Docker.
\end{itemize}
